\documentclass[11pt]{article}
% ======================================================================
%  weekpreamble.tex — shared preamble for the weekly course summaries (EN)
%  Concise, readable, intuition-first. Same box style as the main doc.
%  Compiles with pdflatex.
% ======================================================================
\usepackage[utf8]{inputenc}
\usepackage[T1]{fontenc}
\usepackage[english]{babel}
\usepackage[a4paper,margin=2.2cm]{geometry}
\usepackage{amsmath,amssymb,amsthm,mathtools}
\usepackage{bm}
\usepackage{enumitem}
\usepackage{xcolor}
\usepackage[most]{tcolorbox}
\usepackage{microtype}

\emergencystretch=3em
\allowdisplaybreaks[1]
\setlength{\parindent}{0pt}
\setlength{\parskip}{0.45em}

% ---------- math macros (same names as the main document) ----------
\newcommand{\E}{\mathbb{E}}
\DeclareMathOperator{\Var}{Var}
\DeclareMathOperator{\Cov}{Cov}
\DeclareMathOperator{\Corr}{Corr}
\DeclareMathOperator*{\argmin}{arg\,min}
\DeclareMathOperator{\MSE}{MSE}
\DeclareMathOperator{\trace}{tr}
\newcommand{\Reals}{\mathbb{R}}
\newcommand{\Ints}{\mathbb{Z}}
\newcommand{\Comp}{\mathbb{C}}
\newcommand{\Nat}{\mathbb{N}}
\newcommand{\Prob}{\mathbb{P}}
\newcommand{\Tr}{^{\mathsf{T}}}
\newcommand{\Herm}{^{\mathsf{H}}}
\newcommand{\conj}[1]{\overline{#1}}
\newcommand{\abs}[1]{\left\lvert #1 \right\rvert}
\newcommand{\norm}[1]{\left\lVert #1 \right\rVert}
\newcommand{\ind}{\mathbf{1}}
\newcommand{\eps}{\varepsilon}
\newcommand{\diff}{\nabla}
\newcommand{\acvs}{\gamma}
\newcommand{\acf}{\rho}
\newcommand{\pacf}{\alpha}
\newcommand{\sdf}{\mathcal{S}}
\newcommand{\filt}{\mathcal{L}}
\newcommand{\Bsh}{B}
\newcommand{\ms}{\;\overset{\mathrm{ms}}{=}\;}
\newcommand{\toprob}{\xrightarrow{P}}
\newcommand{\todist}{\xrightarrow{d}}
\newcommand{\iid}{\overset{\text{iid}}{\sim}}

\setlist[itemize]{leftmargin=1.5em,topsep=2pt,itemsep=2pt}
\setlist[enumerate]{leftmargin=1.7em,topsep=2pt,itemsep=2pt}

% ---------- compact, titled (unnumbered) boxes ----------
\tcbset{wkbase/.style={enhanced,breakable,fonttitle=\bfseries,coltitle=white,
  boxrule=0.6pt,arc=1.2mm,left=2.2mm,right=2.2mm,top=1.1mm,bottom=1.1mm,
  before skip=6pt,after skip=6pt,
  attach boxed title to top left={yshift=-3.2mm,xshift=2.4mm},
  boxed title style={boxrule=0pt,arc=0.5mm}}}

\newtcolorbox{defn}[1]{wkbase, colback=blue!4!white, colframe=blue!58!black,
  colbacktitle=blue!58!black, title={Definition --- #1}, top=3mm}
\newtcolorbox{result}[1]{wkbase, colback=red!4!white, colframe=red!60!black,
  colbacktitle=red!60!black, title={#1}, top=3mm}
\newtcolorbox{intu}[1][Intuition]{wkbase, colback=yellow!12!white,
  colframe=orange!72!black, colbacktitle=orange!72!black, coltitle=white,
  title={#1}, top=3mm}
\newtcolorbox{retenir}[1][Key takeaways --- the week's reflexes]{wkbase,
  colback=green!5!white, colframe=green!48!black, colbacktitle=green!48!black,
  title={#1}, top=3mm}
\newtcolorbox{exmpl}[1][Worked example]{wkbase, colback=black!3!white,
  colframe=black!55!white, colbacktitle=black!55!white, title={#1}, top=3mm}

% ---------- week header ----------
\newcommand{\weekheader}[4]{%
  {\small\sffamily\bfseries\color{black!60} TIME SERIES~~$\cdot$~~MATH-342, EPFL}\par
  \vspace{2pt}
  {\LARGE\bfseries Week #1 --- #3\par}
  \vspace{1pt}
  {\itshape\color{black!55} #2}\par
  \vspace{6pt}
  \begin{tcolorbox}[enhanced,colback=blue!3!white,colframe=blue!45!black,
    arc=1.4mm,boxrule=0.7pt,left=3mm,right=3mm,top=2mm,bottom=2mm,before skip=2pt,after skip=10pt]
    \textbf{The big picture.}~ #4
  \end{tcolorbox}
}

\begin{document}
\weekheader{04}{12 March 2026}{Parametric Estimation: Yule--Walker \& Least Squares}{AR/ARMA parameters are estimated by method-of-moments (Yule--Walker), by least squares (forward/backward), or by recursively reconstructing the innovations.}

We have modelled $\{X_t\}$ as an ARMA process; now we \emph{fit} it. Given data $X_1,\dots,X_n$, estimate the AR coefficients $\phi_1,\dots,\phi_p$, the MA coefficients $\theta_1,\dots,\theta_q$, and the innovation variance $\sigma_\eps^2$. Throughout $\eps_t\iid\mathcal N(0,\sigma_\eps^2)$ is Gaussian white noise and the AR part is \textbf{causal} (roots of $\Phi(z)$ outside the unit circle), so that $\eps_t$ is uncorrelated with the past $X_{t-k}$, $k>0$ --- the one fact that makes everything work.

\begin{intu}[The two big ideas]
There are two routes. \textbf{(1) Yule--Walker} = match moments: turn the AR recursion into a linear system in the autocovariances and solve once. \textbf{(2) Least squares} = minimise the squared one-step prediction errors $\sum\eps_t^2$. For pure AR these coincide asymptotically; for ARMA, least squares is the only one that generalises, but the innovations must first be rebuilt.
\end{intu}

\section*{Key definitions}

\begin{defn}{Yule--Walker estimators (AR$(p)$)}
For a mean-zero causal AR$(p)$, form the sample ACVS (mean known to be $0$)
\[
\hat\acvs_\tau=\frac1n\sum_{t=1}^{n-\abs\tau}X_tX_{t+\abs\tau},
\]
collect $\hat{\bm\acvs}_p=(\hat\acvs_1,\dots,\hat\acvs_p)\Tr$ and the symmetric Toeplitz matrix $\hat\Gamma_p=(\hat\acvs_{i-j})_{i,j=1}^p$ (diagonal $\hat\acvs_0$). Then
\[
\hat{\bm\phi}_p=\hat\Gamma_p^{-1}\hat{\bm\acvs}_p,
\qquad
\hat\sigma_\eps^2=\hat\acvs_0-\sum_{j=1}^p\hat\phi_j\,\hat\acvs_j .
\]
This is the empirical analogue of the population identities $\bm\acvs=\Gamma\bm\phi$ and $\acvs_0=\sum_j\phi_j\acvs_j+\sigma_\eps^2$.
\end{defn}

\begin{defn}{Forward least-squares estimator (AR$(p)$)}
Write the AR equations for the observable times $t=p+1,\dots,n$ as $\mathbf X_F=F\bm\phi+\bm\eps$, where $\mathbf X_F=(X_{p+1},\dots,X_n)\Tr$ and the rows of $F$ are the lagged values $(X_{t-1},\dots,X_{t-p})$. The forward LS estimator minimises
\[
\mathrm{SS}_F(\bm\phi)=\norm{\mathbf X_F-F\bm\phi}^2
=\sum_{t=p+1}^n\Big(X_t-\textstyle\sum_{k=1}^p\phi_kX_{t-k}\Big)^2,
\]
giving the OLS solution $\hat{\bm\phi}_F=(F\Tr F)^{-1}F\Tr\mathbf X_F$ and $\hat\sigma_F^2=\norm{\mathbf X_F-F\hat{\bm\phi}_F}^2/(n-2p)$. The first $p$ observations are discarded as conditioning values.
\end{defn}

\begin{defn}{Backward \& forward--backward LS (AR$(p)$)}
By \emph{time reversibility}, regress the \emph{first} $n-p$ values on their \emph{future} lags: $X_t=\sum_{j=1}^p\phi_jX_{t+j}+\nu_t$, $t=1,\dots,n-p$, i.e.\ $\mathbf X_B=B\bm\phi+\bm\nu$, giving $\hat{\bm\phi}_B=(B\Tr B)^{-1}B\Tr\mathbf X_B$ with the same $\bm\phi$. The \textbf{forward--backward} estimator $\hat{\bm\phi}_{FB}$ minimises the combined criterion $\mathrm{SS}_F(\bm\phi)+\mathrm{SS}_B(\bm\phi)$ (stack $F,B$ and $\mathbf X_F,\mathbf X_B$, then OLS). Simulations show it usually beats both.
\end{defn}

\begin{defn}{Least squares for ARMA via residual reconstruction}
For $\Phi(\Bsh)X_t=\Theta(\Bsh)\eps_t$ the innovations are \emph{unobserved}. LS chooses parameters minimising $S(\bm\phi,\bm\theta)=\sum_t\hat\eps_t^{\,2}$, where for each candidate value the residuals are \textbf{reconstructed recursively} from the model equation after setting the $q$ pre-sample innovations (and any pre-sample data) to zero, e.g.\ $\eps_0=\eps_{-1}=\dots=0$. Nonlinear in general; closed form only in the pure-AR case (then it is forward LS).
\end{defn}

\section*{Key results \& formulas}

\begin{result}{Proposition --- Yule--Walker equations}
For a mean-zero causal AR$(p)$ with ACVS $\acvs_\tau$, for every $k\ge1$
\[
\acvs_k=\sum_{j=1}^p\phi_j\,\acvs_{k-j},\qquad\text{i.e.}\qquad
\bm\acvs=\Gamma\bm\phi,\quad \acvs_0=\sum_{j=1}^p\phi_j\acvs_j+\sigma_\eps^2.
\]
\textit{Idea:} multiply $X_t=\sum_j\phi_jX_{t-j}+\eps_t$ by $X_{t-k}$ and take expectations; causality kills $\E[\eps_tX_{t-k}]=0$ for $k>0$. Replacing $\acvs$ by $\hat\acvs$ gives the YW estimators.
\end{result}

\begin{intu}[Why causality is the whole trick]
``Causal $\Rightarrow\E[\eps_tX_{t-k}]=0$'' is the heart of it. $X_{t-k}$ is built from innovations \emph{strictly before} $\eps_t$, and white noise is uncorrelated, so the cross term vanishes --- leaving a clean linear system in the autocovariances. Without causality the noise leaks into the past and YW fails.
\end{intu}

\begin{result}{Proposition --- YW = best linear predictor}
The $\phi_1,\dots,\phi_k$ minimising the one-step MSE $\E\big[(Y_t-\sum_{i=1}^k\phi_iY_{t-i})^2\big]$ satisfy exactly the Yule--Walker equations $\acvs_j=\sum_i\phi_i\acvs_{j-i}$.
\textit{Idea:} the objective is a positive quadratic $\acvs_0-2\sum_i\phi_i\acvs_i+\sum_{i,j}\phi_i\phi_j\acvs_{i-j}$; setting $\partial/\partial\phi_i=0$ gives YW. So the AR coefficients \emph{are} the optimal linear-prediction weights --- which is why YW and LS agree asymptotically.
\end{result}

\begin{result}{Proposition --- time reversibility of a Gaussian AR}
If $\{X_t\}$ is a stationary Gaussian AR$(p)$, then $Y_t=X_{-t}$ is an AR$(p)$ with the \emph{same} coefficients and innovation variance. Equivalently $X_t=\sum_{j=1}^p\phi_jX_{t+j}+\nu_t$.
\textit{Idea:} a Gaussian process is determined by its ACVS, and $\acvs_{-\tau}=\acvs_\tau$; forward and reversed series have identical second-order structure. This is exactly what \emph{legitimises} the backward regression.
\end{result}

\begin{result}{Levinson--Durbin \& a caveat}
\textbf{Levinson--Durbin:} the YW system $\hat\Gamma_p\hat{\bm\phi}_p=\hat{\bm\acvs}_p$ is symmetric Toeplitz; this structure is exploited to solve it recursively in $O(p^2)$ instead of $O(p^3)$, producing as a by-product the PACF $\hat\pacf_k=\hat\phi_{k,k}$ and the error variance at each order.\\
\textbf{Caveat:} the LS estimators $\hat{\bm\phi}_F,\hat{\bm\phi}_B,\hat{\bm\phi}_{FB}$ are \emph{not} constrained to be stationary ($\hat\Phi(z)$ may have roots inside the unit circle) --- a problem for prediction, but fine for spectral estimation. Yule--Walker, by contrast, \emph{always} yields a stationary fit.
\end{result}

\textbf{Recipe --- recursive innovation reconstruction (MA/ARMA).} Solve the model equation for $\eps_t$, fix the $q$ pre-sample innovations (and data) to zero, recurse forward, then minimise $S=\sum_t\hat\eps_t^2$:
\[
\begin{aligned}
\text{MA(1):}\quad &\eps_t=X_t+\theta\eps_{t-1},\\
\text{ARMA(1,1):}\quad &\eps_t=X_t-\phi X_{t-1}+\theta\eps_{t-1}\quad(X_0=\eps_0=0).
\end{aligned}
\]
Each $\hat\eps_t$ is an explicit function of the parameters; because of the $\theta\hat\eps_{t-1}$ feedback, $S$ is generally \emph{not} quadratic and must be minimised numerically.

\begin{intu}[Yule--Walker vs.\ least squares, in practice]
For large $n$ both give nearly the same AR estimates (e.g.\ $\hat\phi\approx\hat\acf_1$ for AR(1), which is also what method-of-moments gives). The differences are at the margins: YW is one linear solve and always stationary; LS extends cleanly to ARMA and to including a mean, but can stray outside the stationary region.
\end{intu}

\section*{A tiny worked example}

\begin{exmpl}[AR(1) Yule--Walker, by hand]
A length-$200$ series gives $\hat\acvs_0=4.0$, $\hat\acvs_1=2.4$. For $p=1$ the YW system is scalar, $\hat\acvs_1=\hat\phi\,\hat\acvs_0$, so
\[
\hat\phi=\frac{\hat\acvs_1}{\hat\acvs_0}=\frac{2.4}{4.0}=0.6=\hat\acf_1,
\qquad
\hat\sigma_\eps^2=\hat\acvs_0-\hat\phi\,\hat\acvs_1=4.0-0.6(2.4)=2.56.
\]
Check: $\hat\acvs_0(1-\hat\phi^2)=4.0(0.64)=2.56$ matches $\acvs_0=\sigma_\eps^2/(1-\phi^2)$. Stationarity: $\Phi(z)=1-0.6z$ has root $z=1.67$, modulus $>1$, so the fit is causal.
\end{exmpl}

\section*{Key takeaways}

\begin{retenir}
\begin{itemize}
\item \textbf{Fit AR$(p)$ by YW:} estimate $\hat\acvs_0,\dots,\hat\acvs_p$, build the Toeplitz $\hat\Gamma_p$ and vector $\hat{\bm\acvs}_p$, solve $\hat{\bm\phi}_p=\hat\Gamma_p^{-1}\hat{\bm\acvs}_p$, then $\hat\sigma_\eps^2=\hat\acvs_0-\sum_j\hat\phi_j\hat\acvs_j$.
\item \textbf{ACF shortcut:} dividing the whole system by $\hat\acvs_0$ replaces covariances by correlations, $\hat{\bm\phi}_p=\bar\Gamma_p^{-1}\hat{\bm\acf}_p$ ($\hat\acvs_0$ cancels; needed only for $\hat\sigma_\eps^2$).
\item \textbf{LS is regression:} forward/backward/forward--backward are all OLS $\hat{\bm\phi}=(A\Tr A)^{-1}A\Tr\mathbf X$. Remember $\hat\sigma^2$ divides by $n-2p$, not $n$.
\item \textbf{MA/ARMA:} the innovations are latent --- \emph{always} reconstruct them recursively (initial $\eps$'s set to $0$) before squaring; expect a numerical minimisation.
\item \textbf{Reflexes:} YW $\Rightarrow$ always stationary; LS $\Rightarrow$ flexible (ARMA, mean) but may be non-stationary; don't put $\acvs_0$ in the RHS vector.
\end{itemize}
\end{retenir}

\end{document}
